\section{Discussion}\label{sec:discussion}

The screen is a first stage. Its intended
pipeline reduces $p$ coordinates to a manageable $d$, after which a joint
method --- the central tail-index subspace of \citet{gardesPodgorny2025}, a
structured tail regression \citep{wangTsai2009,sasakiTaoWang2026}, or
multivariate local estimation --- can resolve interactions and correlated
neighbours among the retained variables. The two stages solve different
problems: rapid elimination of evidently irrelevant coordinates scales to
$p\gg n$, while accurate joint tail inference does not.

Three limitations are intrinsic. First, activity without marginal
detectability: when every fibre of an active coordinate meets the global
argmax set, the score carries no information about it, at any sample size.
Proposition~\ref{prop:geometry} characterizes this failure exactly, and it
is a property of the target, so enlarging $d$ or $n$ cannot repair it;
iterative or grouped extensions would require new identification arguments.
Second, mixture bias: when the projected slowly varying factor decays
logarithmically, condition (E2) restricts the detectable gaps to
$\Delta_{\min}\gg1/\log n$, and the continuous-mixture example of
Section~\ref{sec:estimation} shows the restriction is intrinsic rather
than an artifact of the proofs. Third, rare favourable configurations:
even with a comfortable population gap, the envelope near its maximum is
estimable only from observations whose remaining active coordinates are
simultaneously near their maximizing values, and the mass of such
configurations can decay exponentially with the number of active
coordinates. The weak-signal model M3 in Section~\ref{sec:simulation}
shows the resulting cost.

The main open theoretical question is the gap between the primitive model
(C1)--(C2) and the estimation conditions (E1)--(E3). Projection creates
slow variation that the primitive model does not control, and we do not
know useful primitive conditions implying uniform triangular-array moduli
for all coordinatewise projections. The score functional also invites
variants: a minimum or a low quantile of the profile would emphasize
depressions confined to short intervals, at the price of new uniformity
arguments for empirical profile extrema.

In practice, tuning matters more than the theory prescribes. The tail
fraction trades Hill bias against variance; the bandwidth additionally
determines how often near-maximizing configurations of the other active
coordinates enter a window. We recommend the workflow of
Sections~\ref{sec:simulation} and~\ref{sec:realdata}: run the screen on
the block of nine neighbouring tunings, rank by the aggregated minimum
rank, read the per-setting ranks as the stability diagnostic --- a
variable promoted by a single setting warrants no interpretation, one
holding its rank across the block does --- retain a generous $d$, and
reserve interpretation for variables that survive. Dependence deserves
the same caution: it can help, by concentrating active configurations,
or hurt, by promoting inactive neighbours, and the direction is
design-specific.

Projection turns a tail-index surface into an upper envelope. Variation
of that envelope is a scalable marginal signal for tail-index activity,
and its absence, characterized by the argmax projections, tells the
analyst exactly what a marginal screen can and cannot see.
